\section{Experiments}
\label{sec:experiments}

\subsection{Experiment 1: Ablation Study}
\begin{chinese}
\subsection*{实验 1：消融研究}
\end{chinese}

\textbf{Setup.} We evaluate four conditions (A: baseline, B: full orchestration, C: Pre-only, D: Executor+Post) on SWE-bench Lite (300 tasks) and GAIA (300 tasks). Each task is sent as a \texttt{-pig} model request (Condition B/C/D) or a passthrough request (Condition A) to the Pigs API.
\begin{chinese}
\textbf{设置。} 我们在 SWE-bench Lite（300 项任务）和 GAIA（300 项任务）上评估四个条件（A：基线，B：完整编排，C：仅 Pre，D：Executor+Post）。每个任务作为 \texttt{-pig} 模型请求（条件 B/C/D）或透传请求（条件 A）发送到 Pigs API。
\end{chinese}

\textbf{Procedure.} For each task:
\begin{chinese}
\textbf{流程。} 对每个任务：
\end{chinese}
\begin{enumerate}[nosep]
    \item Send the task prompt to the Pigs API endpoint.
    \begin{chinese}
    将任务提示发送到 Pigs API 端点。
    \end{chinese}
    \item If tool calls are returned, execute them locally and send results back.
    \begin{chinese}
    如果返回工具调用，本地执行后将结果发回。
    \end{chinese}
    \item Repeat until the API returns a final response (no tool calls).
    \begin{chinese}
    重复直到 API 返回最终响应（无工具调用）。
    \end{chinese}
    \item Evaluate the final response against the benchmark's ground truth.
    \begin{chinese}
    根据基准测试的标准答案评估最终响应。
    \end{chinese}
\end{enumerate}

\textbf{Hypothesis.} Full orchestration (B) outperforms baseline (A) on task completion rate, with the Pre phase contributing planning and the Post phase contributing verification.
\begin{chinese}
\textbf{假设。} 完整编排（B）在任务完成率上优于基线（A），Pre 阶段贡献规划能力，Post 阶段贡献验证能力。
\end{chinese}

\subsection{Experiment 2: Control Marker Effectiveness}
\begin{chinese}
\subsection*{实验 2：控制标记有效性}
\end{chinese}

\textbf{Setup.} We compare two configurations on GAIA: (A) with \texttt{PIGFAIL} enabled (failure triggers replan), (B) with \texttt{PIGFAIL} disabled (failure treated as completion).
\begin{chinese}
\textbf{设置。} 我们在 GAIA 上比较两种配置：（A）启用 \texttt{PIGFAIL}（失败触发重规划），（B）禁用 \texttt{PIGFAIL}（失败视为完成）。
\end{chinese}

\textbf{Procedure.} We count: (1) tasks where PIGFAIL was emitted, (2) tasks where replan led to successful completion, (3) tasks that failed despite replan.
\begin{chinese}
\textbf{流程。} 统计：（1）输出 PIGFAIL 的任务数，（2）重规划后成功完成的任务数，（3）尽管重规划仍失败的任务数。
\end{chinese}

\textbf{Hypothesis.} PIGFAIL-based replan improves outcomes on tasks where the initial approach was flawed.
\begin{chinese}
\textbf{假设。} PIGFAIL 重规划在初始方法有缺陷的任务上改善结果。
\end{chinese}

\subsection{Experiment 3: Three-Protocol Consistency}
\begin{chinese}
\subsection*{实验 3：三协议一致性}
\end{chinese}

\textbf{Setup.} We send the same set of 100 tasks through three API protocols: OpenAI Chat (\texttt{/chat/completions}), Anthropic Messages (\texttt{/v1/messages}), and OpenAI Responses (\texttt{/responses}).
\begin{chinese}
\textbf{设置。} 我们通过三种 API 协议发送相同的 100 项任务：OpenAI Chat（\texttt{/chat/completions}）、Anthropic Messages（\texttt{/v1/messages}）和 OpenAI Responses（\texttt{/responses}）。
\end{chinese}

\textbf{Procedure.} For each task and protocol, we record: (1) whether the task completed, (2) the visible text output, (3) any tool calls made. We measure consistency as the percentage of tasks where all three protocols agree on completion status.
\begin{chinese}
\textbf{流程。} 对每个任务和协议，记录：（1）任务是否完成，（2）可见文本输出，（3）工具调用。一致性定义为三种协议在完成状态上达成一致的任务百分比。
\end{chinese}

\textbf{Hypothesis.} The protocol-native request preservation ensures consistent behavior across all three protocols.
\begin{chinese}
\textbf{假设。} 协议原生请求保真确保三种协议间行为一致。
\end{chinese}

\subsection{Experiment 4: Cost-Benefit Analysis}
\begin{chinese}
\subsection*{实验 4：成本效益分析}
\end{chinese}

\textbf{Setup.} On SWE-bench Lite, we measure per-task: (1) total input tokens, (2) total output tokens, (3) number of API calls, (4) wall-clock latency, (5) task completion.
\begin{chinese}
\textbf{设置。} 在 SWE-bench Lite 上，测量每个任务的：（1）总输入 token，（2）总输出 token，（3）API 调用次数，（4）实际延迟，（5）任务完成。
\end{chinese}

\textbf{Procedure.} We compare Condition A (baseline) vs.\ Condition B (full orchestration) and compute the overhead ratio for each metric.
\begin{chinese}
\textbf{流程。} 对比条件 A（基线）和条件 B（完整编排），计算每个指标的开销比。
\end{chinese}

\textbf{Hypothesis.} The token overhead of orchestration (3 phases vs.\ 1) is justified by improved completion rate.
\begin{chinese}
\textbf{假设。} 编排的 token 开销（3 阶段 vs 1 阶段）通过提升的完成率得到补偿。
\end{chinese}
